\scanpage{373}

We shall thus have semidirect product structures
\[
(16)\qquad
\mathcal G_{\rho,\sigma}\simeq
\mathcal G_{\rho,\sigma}^{\unsure{\flat|b}}(0)\cdot\Lambda_0,
\qquad
\mathcal G_{\rho,\sigma}^{*}\simeq
\mathcal G_{\rho,\sigma}^{*\unsure{\flat|b}}(0)\cdot\Lambda_0,
\]
where
\[
(17)\qquad
\left\{\begin{array}{l}
\displaystyle\Lambda_0\subset\mathcal G_{\rho,\sigma},\quad
\Lambda_0=\left\{\left(
\alpha=\underset{\substack{\Vert\\\varepsilon_0^{\unsure{n|u}}}}{\lambda},
\beta=\underset{\substack{\Vert\\\varepsilon_0^{\unsure{n|u}}}}{\lambda},
\mu=1\right)\,\middle|\,\lambda\in\Gamma L_0\right\}\\
\hfill(\unsure{n|u}\in\Gamma\mathbb \mathbb G_a)\\[4pt]
\text{corresponding to:}\quad L_0\times\{1\}\times\{1\}
\subset B'_0\times T_\rho\times T_\sigma.
\end{array}\right.
\]

We shall finally consider another subgroup \inserted{scheme} of
$\mathcal G_{\rho,\sigma}^{*}$, by the conditions
\[
(18)\qquad
\left\{\begin{array}{l}
\det\alpha\in\unsure{\Gamma\mu_2}\quad\text{i.e. }(\det\alpha)^2=1\\
\det\beta=1
\end{array}\right.
\qquad\left|\quad\text{whence }S\mathcal G_{\rho,\sigma}^{\unsure{*|\flat}}
\subset\mathcal G_{\rho,\sigma}^{*}\right.
\]

\deleted{Let then} \inserted{\unsure{Remark}}
\deleted{$\unsure{\varepsilon'\longleftarrow(\alpha,\beta),|\varepsilon'(\alpha,\beta),}$}
$S\mathcal G_{\rho,\sigma}^{*}$ can be interpreted as the scheme of systems
\[
(\alpha,\beta,\mu,\varepsilon,\varepsilon')
\]
satisfying
\[
(19)\qquad [\alpha,\rho]=\varepsilon[\beta,\sigma]\varepsilon_1(\mu-1),
\quad\det\alpha=\varepsilon',\quad
\inserted{\det\beta=1,}\quad
\underline{\varepsilon^2=\varepsilon'^2=1}
\]
\margin{N.B. $\varepsilon^2=1$ already follows from the first relation.}

In terms of the systems
\[
(\underline u,\underline a,\underline b)\in\Gamma B'_0\times T_\rho\times N'_\sigma
\]
the equations (19) (of which the \unsure{first|first} is automatic), reduce to:
\[
(20)\qquad\det\underline u=\det\underline b=\varepsilon'\det\underline a,
\qquad\varepsilon'^2=1,
\]

\mspage{372}

% Printed archive footer outside the manuscript frame:
% 372
% https://grothendieck.umontpellier.fr
% Université de Montpellier, UMR CNRS 5149, Case courrier 051, Place Eugène Bataillon
% 34095 Montpellier cedex 5 - France

\scanpage{374}
\mspage{488}
which \unsure{show} that $S\mathcal{G}_{\rho,\sigma}^{*}$ is \inserted{indeed} a \uline{subgroup}
of $\mathcal{G}_{\rho,\sigma}^{*}$. This subgroup contains $\Lambda_0$
(\unsure{corresponding}: $L_0\times\unsure{\{1\}\times\{1\}|\{(1,1)\}}\subset B'_0\times T_\rho\times N'_\sigma$) --- which
corresponds to the case where \deleted{\unsure{we have}}
\[
\det\underline{u}=\det\underline{b}=\det\underline{a}=1\qquad(\varepsilon'=1),
\tag{21}
\]
and we have a $\unsure{\frac12}$ direct product structure
\[
S\mathcal{G}_{\rho,\sigma}^{*}\simeq S\mathcal{G}_{\rho,\sigma}^{*\unsure{\flat|b}}(0)\mathbin{.}\Lambda_0
\tag{22}
\]
where
\[
S\mathcal{G}_{\rho,\sigma}^{*}(0)=S\mathcal{G}_{\rho,\sigma}^{*}\cap\mathcal{G}_{\rho,\sigma}^{*\unsure{\flat|b}}(0)
\tag{23}
\]
corresponds to the \unsure{subgroup|subgroup} \deleted{of} $\deleted{\unsure{B'_0}}\unsure{T_0|T'_0}\times T_\rho\times N'_\sigma$
\deleted{of $B'_0\times T_\rho\times N_\sigma$} defined by the same equations
(20). Of course, we shall also have
\[
S\mathcal{G}_{\rho,\sigma},\quad S\mathcal{G}_{\rho,\sigma}^{\unsure{\flat|b}}(0)\quad(\text{without }*)
\tag{24}
\]
\deleted{as} \inserted{intersections} of $S\mathcal{G}_{\rho,\sigma}^{*}$, $S\mathcal{G}_{\rho,\sigma}^{\unsure{\flat|b}*}(0)$ with
$S\mathcal{G}_{\rho\sigma}$ --- deduced from the expression
of these in terms of $(B'_0,\unsure{T_0|T'_0}),T_\rho,N'_\sigma$
by replacing $N'_\sigma$ with $T_\sigma$.

To make explicit the structure of the $S\mathcal{G}_{\rho,\sigma}^{*}$,
we are reduced via (22) \uline{to studying the}
group
\fbox{\diagram{\begin{tikzcd}[ampersand replacement=\&]
(25)\quad S\mathcal{G}_{\rho,\sigma}^{*\unsure{\flat|b}}(0) \arrow[r, hook] \& \unsure{T_0|T'_0}\times T_\rho\times N'_\sigma
\end{tikzcd}}}
\mspage{373}

% Printed archive footer outside the manuscript frame:
% 373
% https://grothendieck.umontpellier.fr
% Université de Montpellier, UMR CNRS 5149, Case courrier 051, Place Eugène Bataillon
% 34095 Montpellier cedex 5 - France

\scanpage{375}
described by equations (20). But we have
\diagram{\begin{tikzcd}[ampersand replacement=\&]
T_0 \arrow[r,"\sim"] \& \mathbb{G}_m \\
\underline{u} \arrow[r,mapsto] \& \det\underline{u}
\end{tikzcd}},
which implies that the projection
\diagram{\begin{tikzcd}[ampersand replacement=\&,column sep=large]
S\mathcal{G}_{\rho,\sigma}^{*\unsure{\flat|b}}(0) \arrow[r] \& T_\rho\times N'_\sigma \\
(\underline{u}\underline{a}^{-1},\underline{u}\underline{b}^{-1},\det\underline{u},\varepsilon=\varepsilon(\underline{b})) \arrow[r] \arrow[d,phantom,"\wr" description] \& (\underline{a},\underline{b}) \\
(\underline{u},\underline{a},\underline{b}) \&
\end{tikzcd}}
is a monomorphism, and induces an isomorphism of
$S\mathcal{G}_{\rho,\sigma}^{*}(0)$ onto the \unsure{subgroup|subgroup} of $T_\rho\times N'_\sigma$
\unsure{defined|defined} by the equation
\[
\det\underline{b}=\varepsilon'\det\underline{a},\qquad \varepsilon'{}^2=1
\tag{21}
\]
we shall have
\[
(\det\underline{b})^2=(\det\underline{a})^2.
\]

In terms of the homomorphisms (\unsure{determinants}, in short)
\diagram{\begin{tikzcd}[ampersand replacement=\&]
T_\rho \arrow[r,"\delta_\rho"] \& \mathbb{G}_m, \& N'_\sigma \arrow[r,"\unsure{\delta_\sigma}"] \& \mathbb{G}_m
\end{tikzcd}}
this amounts to taking the inverse image of the
diagonal \inserted{of $\mathbb{G}_m\times \mathbb{G}_m$} under $\delta_\rho^2,\delta_\sigma^2$.

We see at once (since $\delta_\rho,\delta_\sigma$
are epimorphisms) that the homomorphism
\diagram{\begin{tikzcd}[ampersand replacement=\&,column sep=large]
(22)\quad S\mathcal{G}_{\rho,\sigma}^{*\unsure{\flat|b}}(0) \arrow[r] \& \mu_2\times\mu_2 \\
(\alpha,\beta,\mu,\varepsilon,\varepsilon') \arrow[r,mapsto] \& (\varepsilon,\varepsilon') \\
\text{i.e.}\quad(\underline{u},\underline{a},\underline{b}) \arrow[r,mapsto] \& \left(\underline{\varepsilon}(\underline{b}),\underbrace{\det\underline{b}/\det\underline{a}}_{\det(\underline{b}\underline{a}^{-1})}\right)
\end{tikzcd}}
is an epimorphism, and \deleted{the kernel} thus gives
rise to an exact sequence
\mspage{374}

% Printed archive footer outside the manuscript frame:
% 374
% https://grothendieck.umontpellier.fr
% Université de Montpellier, UMR CNRS 5149, Case courrier 051, Place Eugène Bataillon
% 34095 Montpellier cedex 5 - France

\scanpage{376}
\mspage{489}
\diagram{\begin{tikzcd}[ampersand replacement=\&,column sep=small]
(23)\quad1 \arrow[r] \& S\mathcal{G}_{\rho,\sigma}^{*\circ\unsure{\flat|b}}(0) \arrow[r] \& S\mathcal{G}_{\rho,\sigma}^{*\unsure{\flat|b}}(0) \arrow[r] \& \mu_2\times\mu_2 \arrow[r] \& 1
\end{tikzcd}}
where $S\mathcal{G}_{\rho,\sigma}^{\unsure{\circ\flat|\circ}}(0)$ is isomorphic to the subscheme of
$T_\rho\times T_\sigma$ defined by the equation
\[
\det\underline{a}=\det\underline{b}
\tag{24}
\]
and thus in turn breaks down into
\diagram{\begin{tikzcd}[ampersand replacement=\&,column sep=small]
(25)\quad1 \arrow[r] \& ST_\rho^{\unsure{\flat|b}}\unsure{\times} ST_\sigma^{\unsure{\flat|b}} \arrow[r] \& S\mathcal{G}_{\rho,\sigma}^{\circ\unsure{\flat|b}}(0) \arrow[r] \& \mathbb{G}_m \arrow[r] \& 1 \\
\& \& \deleted{\unsure{(\underline{u},|\underline{u}}}(\underline{a},\underline{b}) \arrow[r,mapsto] \& \deleted{\det\underline{u}}\det\underline{a}=\det \&
\end{tikzcd}}
where we have set
\[
\left\{
\begin{aligned}
ST_\rho^{\unsure{\flat|b}}&=\operatorname{Ker}\left(T_\rho\xrightarrow{\delta_\rho}\mathbb{G}_m\right)\\
ST_\sigma^{\unsure{\flat|b}}&=\operatorname{Ker}\left(T_\sigma\xrightarrow{\delta_\sigma}\mathbb{G}_m\right)
\end{aligned}
\right.
\tag{26}
\]

If for example $\rho$ \inserted{(\unsure{non-s|non-g})} is $\frac12$ simple, hence $T_\rho$ a
maximal torus, then $ST_\rho^{\unsure{\flat|b}}$ is itself a torus \inserted{(of \unsure{dimension 1})} ---
\unsure{ditto} for $\sigma$. Thus, \deleted{\illegible{1}} in the ``modular case''
\[
\rho=\begin{pmatrix}0&1\\-1&1\end{pmatrix},\qquad
\sigma=\begin{pmatrix}0&-1\\1&0\end{pmatrix},
\]
$ST_\rho^{\unsure{\flat|b}}$ is a
torus \deleted{of} except in characteristic 3, and $ST_\sigma^{\unsure{\flat|b}}$ a torus except in
characteristic 2. \uline{Thus}, \uline{except in characteristic \inserted{2,3}}, $S\mathcal{G}_{\rho,\sigma}^{\circ\unsure{\flat|b}}(0)$ is a
\uline{torus of dimension 3}.

To finish determining the structure
of $S\mathcal{G}_{\rho,\sigma}^{*}$ in terms of $\mathcal{G}_{\rho,\sigma}^{*\unsure{\circ}}(0)$ and
of the decomposition (22) as a semidirect product, we must \deleted{specify}\inserted{make explicit} the operation of
\mspage{375}

% Printed archive footer outside the manuscript frame:
% 375
% https://grothendieck.umontpellier.fr
% Université de Montpellier, UMR CNRS 5149, Case courrier 051, Place Eugène Bataillon
% 34095 Montpellier cedex 5 - France

\scanpage{377}
$S\mathcal{G}_{\rho,\sigma}^{*\circ}(0)$ on $\Lambda_0$, by \unsure{interpreting}
$S\mathcal{G}_{\rho,\sigma}^{*\circ}(0)$ as a subgroup of $T_\rho\times N'_\sigma$,
and $\Lambda_0$ as \inserted{$L_0$} $\simeq \mathbb{G}_{\unsure{a|0}}$ via $\lambda\mapsto\varepsilon_0(\lambda)$.

Now the relations
\[
\deleted{\underline{u}\unsure{\varepsilon_0(\lambda)}\underline{u}^{-1}=}
\]
\[
\underline{u}(\varepsilon_0(\lambda))\underline{u}^{-1}
 =\varepsilon_0(\det(\underline{u})\lambda),\qquad
\det(\underline{u})=\det(\underline{b})=\delta_\sigma(\underline{b})
\]
show that the operation takes place via
the composite character
\diagram{\begin{tikzcd}[ampersand replacement=\&]
(27)\quad S\mathcal{G}_{\rho,\sigma}^{*\unsure{\circ|\flat}}(0) \arrow[r] \& N'_\sigma \arrow[r,"\delta_\sigma"] \& \mathbb{G}_m
\end{tikzcd}}, OK.

For what follows, we shall need
a canonical homomorphism
\diagram{\begin{tikzcd}[ampersand replacement=\&,column sep=large,row sep=large]
(28)\quad\mathcal{G}_{\rho,\sigma}^{*} \arrow[rr,"\unsure{\theta|\Theta}"] \arrow[dr,shift right=0.5ex,"{\deleted{\begin{array}{cccc}(\unsure{\alpha},&\beta,&\mu,&\unsure{\varepsilon})\\\unsure{\Vert|=}&\unsure{\Vert|=}&&\\\underline{u}\underline{a}^{-1}&\underline{u}\underline{b}^{-1}&&\end{array}}}"'] \& \& G \\
\& B'_0\times T_\rho\times\unsure{N'_\sigma|T'_\sigma} \arrow[ul,shift right=0.5ex] \arrow[ur,dashed] \&
\end{tikzcd}}
which identifies with the \uline{first} \uline{projection}
\diagram{\begin{tikzcd}[ampersand replacement=\&]
(\underline{u},\underline{a},\underline{b}) \arrow[r,mapsto] \& \underline{u}
\end{tikzcd}}
of the product, but \unsure{considered|considered} as
a morphism into $G$ thanks to the inclusion
\diagram{\begin{tikzcd}[ampersand replacement=\&]
B'_0 \arrow[r,phantom,"{\unsure{\hookrightarrow|\subset}}" description] \& G.
\end{tikzcd}}
We shall also \unsure{often} set
\unsure{(}
\[
\underline{u}=\unsure{\theta|\Theta}(\alpha,\beta,\mu,\varepsilon),
\]
\unsure{thus|whence} \deleted{if $\det$} we shall have
\mspage{376}

% Printed archive footer outside the manuscript frame:
% 376
% https://grothendieck.umontpellier.fr
% Université de Montpellier, UMR CNRS 5149, Case courrier 051, Place Eugène Bataillon
% 34095 Montpellier cedex 5 - France